\begin{center}
    \Large{\textbf{TÓM TẮT NỘI DUNG ĐỒ ÁN}}\\
\end{center}
\vspace{1cm}

Trong các hệ thống blockchain công khai như Ethereum, mọi giao dịch đều được ghi nhận minh bạch và có thể truy xuất công khai. Mặc dù đặc điểm này đảm bảo tính minh bạch và khả năng xác minh của mạng lưới, nó đồng thời bộc lộ hạn chế về quyền riêng tư của người dùng khi toàn bộ lịch sử giao dịch, số dư tài sản và mối liên hệ giữa các địa chỉ ví đều có thể bị truy vết. Hiện nay, một số giải pháp như mixer hoặc các blockchain tập trung vào quyền riêng tư đã được áp dụng nhằm che giấu thông tin giao dịch; tuy nhiên, các giải pháp này thường tồn tại những rào cản về khả năng tích hợp, chi phí triển khai, hoặc mức độ tương thích với hệ sinh thái Ethereum. Trong bối cảnh đó, cơ chế Stealth Address được đánh giá là một giải pháp kỹ thuật khả thi nhằm gia tăng tính ẩn danh cho người nhận thông qua việc tạo ra một địa chỉ nhận duy nhất cho mỗi giao dịch.

Đồ án này lựa chọn hướng tiếp cận kết hợp Stealth Address với Account Abstraction trên Ethereum nhằm tận dụng khả năng tùy biến của smart wallet đồng thời nâng cao quyền riêng tư trong giao dịch tài sản số. Trên cơ sở đó, hệ thống được thiết kế cho phép người nhận đăng ký khóa công khai, người gửi tạo địa chỉ stealth tương ứng cho từng giao dịch và người nhận có thể phát hiện, quản lý tài sản thông qua ví thông minh.

Đóng góp chính của đồ án là đề xuất và xây dựng mô hình tích hợp Stealth Address với Account Abstraction trên nền tảng Ethereum, qua đó hỗ trợ giao dịch riêng tư hơn mà vẫn tương thích với hạ tầng blockchain hiện có. Kết quả đạt được là một hệ thống thử nghiệm chứng minh khả năng áp dụng cơ chế Stealth Address trong việc bảo vệ quyền riêng tư của người dùng trên mạng Ethereum.

\begin{flushright}
Sinh viên thực hiện\\
\begin{tabular}{@{}c@{}}
\textit{(Ký và ghi rõ họ tên)}
\end{tabular}
\end{flushright}